\chapter{距離空間と測度の準備}
\label{ch:found-preliminaries}

本編で用いる実数・距離空間・測度論の準備をまとめる．

\section{実数の上限・下限}

\begin{definition}{上限と下限}{meas-sup-inf}
  $A\subset\R$ を空でない集合とする．
  \begin{itemize}
    \item $A$ が上に有界であるとき，$A$ の\textbf{上限}（supremum）を
      \[
        \sup A := \min\{M\in\R\mid a\leq M\;\;(\forall a\in A)\}
      \]
      と定義する．すなわち $\sup A$ は $A$ の上界のうち最小のものである．
      $A$ が上に有界でないとき $\sup A := +\infty$ とする．
    \item $A$ が下に有界であるとき，$A$ の\textbf{下限}（infimum）を
      \[
        \inf A := \max\{m\in\R\mid m\leq a\;\;(\forall a\in A)\}
      \]
      と定義する．すなわち $\inf A$ は $A$ の下界のうち最大のものである．
      $A$ が下に有界でないとき $\inf A := -\infty$ とする．
  \end{itemize}
\end{definition}

\section{測度の台と指示関数}

\begin{definition}{測度の台}{meas-support}
  確率測度 $\mu$ の\textbf{台}とは
  \[
    \supp\mu
    =
    \left\{x\in X\mid x\text{ の任意の開近傍 }U\text{ に対して }\mu(U)>0\right\}
  \]
  のことである．
\end{definition}

\begin{definition}{指示関数}{meas-indicator}
  集合 $A\subset X$ の\textbf{指示関数} $\mathbf{1}_A\colon X\to\{0,1\}$ を
  \[
    \mathbf{1}_A(x):=
    \begin{cases}
      1 & (x\in A),\\
      0 & (x\notin A)
    \end{cases}
  \]
  と定める．$\mu$ が確率測度のとき
  $\int_X\mathbf{1}_A\,\d\mu=\mu(A)$ が成り立つ．
\end{definition}

\section{収束定理}

\begin{definition}{各点収束}{meas-pointwise-convergence}
  関数列 $f_n\colon X\to\R$ が $f\colon X\to\R$ に\textbf{各点収束}するとは，
  すべての $x\in X$ に対して
  \[
    \lim_{n\to\infty}f_n(x)=f(x)
  \]
  が成り立つことをいう．
  さらに各 $x$ で $f_1(x)\geq f_2(x)\geq\cdots$ であるとき
  \textbf{単調減少に各点収束する}といい，$f_n\downarrow f$ と書く．
\end{definition}

\begin{theorem}{有界収束定理}{meas-bounded-convergence}
  $(X,\Borel(X),\mu)$ を有限測度空間とする．
  関数列 $f_n\colon X\to\R$ が $f$ に各点収束し，
  ある定数 $M<\infty$ で $\abs{f_n(x)}\leq M$（$\forall n,\,\forall x$）が
  成り立つならば，
  \[
    \lim_{n\to\infty}\int_X f_n\,\d\mu=\int_X f\,\d\mu.
  \]
\end{theorem}

\begin{proof}
  \[
    g_n(x) := \sup_{k \geq n} \abs{f_k(x) - f(x)}
  \]
  とおく．$g_n$ は非負可測であり，$\abs{f_n - f} \leq g_n$ をみたす．
  また $g_n$ は単調減少（$g_1 \geq g_2 \geq \cdots$）であり，
  $f_n \to f$ の各点収束から $g_n(x) \to 0$（すべての $x$）である．
  さらに $\abs{f_k} \leq M$ と $\abs{f} \leq M$ より $g_n \leq 2M$ である．

  $h_n := 2M - g_n$ とおくと，$h_n$ は非負かつ単調増加で
  $h_n \uparrow 2M$ である．
  単調収束定理より
  \[
    \int_X h_n \, \d\mu
    \;\uparrow\;
    \int_X 2M \, \d\mu
    = 2M \, \mu(X).
  \]
  $\mu(X) < \infty$ だから
  \[
    \int_X g_n \, \d\mu
    = 2M \, \mu(X) - \int_X h_n \, \d\mu
    \;\to\; 0.
  \]
  $\abs{f_n - f} \leq g_n$ より
  \[
    \left\lvert\int_X f_n \, \d\mu - \int_X f \, \d\mu\right\rvert
    \leq \int_X \abs{f_n - f} \, \d\mu
    \leq \int_X g_n \, \d\mu
    \to 0.  \qedhere
  \]
\end{proof}

\section{$L^p$ ノルムと積分不等式}

\begin{definition}{$L^p$ ノルム}{meas-lp-norm}
  測度空間 $(\Omega,\mathcal{A},\mu)$ 上の可測関数 $f\colon\Omega\to\R$ に対し，
  $1\leq p<\infty$ のとき
  \[
    \norm{f}_{L^p(\mu)}
    :=
    \left(\int_\Omega\abs{f}^p\,\d\mu\right)^{1/p},
  \]
  $p=\infty$ のとき
  \[
    \norm{f}_{L^\infty(\mu)}
    :=
    \inf\{M\geq 0\mid \abs{f}\leq M\;\;\mu\text{-a.e.}\}
  \]
  と定める．
\end{definition}
